\chapter{Conclusion and Further Work}


\section{Introduction}


This chapter concludes the thesis by summarizing the research contributions, objective achievement, quantitative results, limitations, challenges, and future work. The conclusion is written in relation to the research aim: to develop and evaluate an explainable multimodal fake news detection framework that adaptively fuses image and text predictions using reinforcement learning.


\section{Summary of Research Contributions}


The main contribution of this research is a transparent reliability-aware adaptive fusion framework for multimodal fake news detection. The framework combines a ResNet-based image branch, a DistilBERT-based text branch, reliability feature extraction, reinforcement learning-based fusion, baseline evaluation, robustness analysis, and explanation generation.


The first contribution is the formulation of image-text fusion as a reliability-aware contextual-bandit-style action-selection problem. Instead of using a fixed fusion rule for all samples, the system observes a state containing prediction, confidence, quality, and disagreement features, then selects an interpretable fusion action.


The second contribution is the implementation of a reproducible experimental pipeline. The project includes dataset preparation, image downloading, model training, inference, reliability computation, baseline comparison, RL fusion, supervised fusion, ablation, robustness, threshold tuning, seed checks, and explainability outputs.


The third contribution is the explanation mechanism. The framework produces Grad-CAM image explanations, token-level text saliency, and fusion-level action/weight explanations.


\section{Achievement of Objectives}


The first objective was to conduct a critical literature review. This was addressed by reviewing fake news detection, multimodal learning, fusion strategies, reinforcement learning, and explainable AI. The review identified the gap for sample-level adaptive and explainable fusion.


The second objective was to study relevant technologies and theoretical foundations. This was addressed through the study of transformer-based text modelling, residual image modelling, reinforcement learning, reliability estimation, and explanation methods.


The third objective was to design a reliability-aware adaptive fusion approach. This was achieved by defining a state representation, action space, reward design, and fusion workflow.


The fourth objective was to implement a reproducible framework. This was achieved through the project repository, notebooks, configuration files, source code modules, and output evidence files.


The fifth objective was to evaluate the framework. This was achieved through held-out test evaluation, baseline comparison, ablation, robustness analysis, explainability examples, and seed stability checks.


\section{Quantitative Performance Summary}


The final evaluation used 26,471 available multimodal samples from the Fakeddit-derived final subset. The image-only branch achieved 0.7769 test macro F1, while the text-only branch achieved 0.8388. The best fixed fusion baselines achieved approximately 0.8673 test macro F1. The main RL adaptive fusion run achieved 0.8676 test macro F1, and validation-selected threshold tuning improved it to 0.8685.


The ablation study showed that the full reliability-aware state improved over reduced state representations. Probabilities alone achieved 0.8523 test macro F1, prediction and confidence achieved 0.8590, and the full reliability state achieved 0.8673. This supports the usefulness of reliability-aware state features.


The matched supervised MLP fusion baseline achieved 0.8547 macro F1 at the default threshold and 0.8565 after threshold tuning. The RL fusion result was higher in the main final run, and paired McNemar testing showed a significant difference against the matched supervised MLP controller (p=0.000045). However, multi-seed checks showed that the RL improvement over fixed fusion was not stable across all tested seeds. Therefore, the performance conclusion is stated cautiously.


\section{Limitations of the Study}


The first limitation is the marginal improvement over strong fixed fusion. The selected RL run slightly exceeded equal fusion, but paired testing showed that this difference was not statistically significant. This affects the interpretation of the contribution: the work should be defended as transparent adaptive fusion with competitive performance, not as a clearly superior accuracy model.


The second limitation is seed stability. The main RL run achieved the strongest reported result, but additional seeds showed weaker RL performance against equal fusion. This indicates that the controller is sensitive to training conditions and requires stronger tuning or stability mechanisms before deployment.


The third limitation is dataset scope and external validation. The final evaluation used a subset derived from one dataset source, so the main results demonstrate within-dataset generalization. The OpenFake case study showed that cross-dataset transfer was difficult: the text branch retained ranking signal after calibration, but the image branch and Fakeddit-trained fusion policy did not transfer strongly in zero-shot form. This limits claims about real-world generalization.


The fourth limitation is image availability filtering. The final multimodal subset excluded posts whose images could not be downloaded. This was necessary for image-text evaluation, but it may introduce availability bias because unavailable images may not be random.


The fifth limitation is heuristic reliability measurement. Image blur, brightness, contrast, and entropy measure visual usability rather than truthfulness. Text length, availability, and truncation measure input usability rather than semantic reliability. These proxies helped structure the fusion state, but future work should consider calibrated uncertainty, predictive entropy, perplexity, Monte Carlo dropout, or learned reliability estimators.


The sixth limitation is text-branch overfitting. The text branch showed a large train-test performance gap. Although validation-based checkpoint selection was used, stronger regularization should be investigated.


The seventh limitation is the discrete action space and binary reward. The action set supports interpretability, but it restricts fusion to a small number of weight pairs. The binary reward is clear but does not reward calibration or probability margin.


The eighth limitation is explainability validation. A deletion-based faithfulness check was conducted on 300 samples, but the results were mixed: text token saliency showed positive faithfulness, while image Grad-CAM did not outperform random image-region deletion in this metric. Therefore, image explanations should be treated as qualitative visual evidence rather than fully validated causal explanations.


The ninth limitation is that several analyses were grouped or subset-based rather than exhaustive. Grouped ablations, robustness tests, and faithfulness checks support interpretation, but they do not replace larger statistical validation across datasets and seeds.


\section{Challenges Encountered}


Several practical challenges were encountered during the research. Image downloading from the dataset was incomplete because some external image URLs were unavailable. Computational limitations required the use of manageable final subsets and GPU runtime planning. The fusion controller required careful baseline comparison because small numerical improvements can be misleading if training budgets are not controlled. Finally, explainability outputs had to be designed so that image, text, and fusion-level evidence could be inspected together.


\section{Future Work}


Future work should evaluate the framework on additional multimodal fake news datasets and cross-domain event splits. The OpenFake case study indicates that future versions should include external-domain calibration and domain adaptation rather than assuming that probabilities learned from one dataset will transfer unchanged. The RL fusion controller should also be trained with more seeds, tuned more systematically, and compared using statistical tests on larger evaluation protocols.


The text branch could be improved using stronger regularization, validation-loss-aware checkpoint selection, or domain-specific pretraining. The reward function could be extended from binary correctness to margin-based or log-loss-based rewards. The action space could also be expanded or converted into a continuous policy if higher fusion precision is required.


Finally, explainability evaluation should be expanded beyond the current deletion-based check using insertion tests, larger perturbation protocols, and human-centered studies to determine whether the generated explanations are useful and trustworthy.


\section{Summary}


This thesis presented a reinforcement learning-based adaptive fusion framework for multimodal fake news detection. The framework demonstrates that reliability-aware adaptive fusion can provide competitive multimodal performance while exposing interpretable sample-level modality decisions through action records, modality weights, and branch-level evidence. The contribution is therefore a transparent and reproducible adaptive fusion framework rather than a claim of superior predictive accuracy. Future work should improve policy stability, uncertainty estimation, external validation, domain adaptation, and explanation faithfulness evaluation.

\clearpage
